\chapter{Probability I}\label{ch:iv:probability-i}

``I roll two dice and tell you that at least one of them is a six. What is the chance that both are?'' Half the
room says one sixth, a few say one eleventh, and one candidate asks how the interviewer came to say it, which
is the only question that settles the answer. Probability questions in interviews are short, and most of them
test one of four things: counting without double counting, conditioning on the right event, seeing a symmetry,
and knowing that what you are told depends on how you came to be told it.

\section{Counting}

Every finite probability is a ratio of counts over equally likely outcomes, and the difficulty is always in the
count. Four tools cover most interview questions: the product rule (choices made in sequence multiply);
arrangements of a word with repeated letters, $n!/(k_1!\,k_2!\cdots)$; subsets, $\binom nk$; and multisets, the
number of ways to put $n$ identical items in $k$ labelled boxes, $\binom{n+k-1}{k-1}$ (``stars and bars''), or
$\binom{n-1}{k-1}$ if every box must get at least one.

\begin{definition}[Complementary counting]\label{def:iv:probability-i:complement}
\emph{Complementary counting}\index{complementary counting} computes the probability of an event through its
complement, $\P(A) = 1 - \P(A^c)$, when the complement is a single simple case and the event itself is a union of
many overlapping ones: ``at least one'', ``some two coincide'', ``not all different''.
\end{definition}

\begin{example}[Colliding identifiers]\label{ex:iv:probability-i:ids}
Thirty client order identifiers are drawn uniformly and independently from 10\,000 values. The chance that all
differ is $\prod_{k=0}^{29}(1 - k/10\,000) \approx 0.957$, so at least two coincide with probability about 0.043.
The approximation $1 - e^{-n(n-1)/(2N)}$ gives the same to three decimals, and it shows the scale: collisions
become likely once $n$ is of the order of $\sqrt N$.
\end{example}

\section{Conditioning and Bayes}

Conditioning restricts the sample space to the outcomes consistent with what is known and renormalises. Bayes's
rule is the same statement written in the direction the question needs. For binary hypotheses it is fastest in
odds: posterior odds equal prior odds times the likelihood ratio $\P(\text{data} \mid H)/\P(\text{data} \mid
H^c)$, and independent pieces of evidence multiply their likelihood ratios (One Quant Book~4, chapter~14, on
priors and posteriors).

\begin{method}[Three ways to condition]\label{met:iv:probability-i:condition}
\begin{enumerate}
\item \emph{Tree}: draw the stages in the order they happen, write the probabilities on the branches, and add the
leaves consistent with what is known.
\item \emph{Table}: when two binary variables are involved, a two-by-two table of counts out of 1\,000 or
10\,000 makes base rates visible.
\item \emph{Odds}: prior odds times likelihood ratio, when the evidence arrives in pieces.
\end{enumerate}
\end{method}

\begin{example}[A rare event and a good test]\label{ex:iv:probability-i:bayes}
One account in 500 is fraudulent; an alert fires on 95\% of fraudulent accounts and on 2\% of the others. Out of
100\,000 accounts, 200 are fraudulent and 190 of them trigger the alert; of the 99\,800 others, 1\,996 do. An
alert therefore means fraud with probability $190/2\,186 \approx 0.087$. The prior odds of $1:499$ times the
likelihood ratio $0.95/0.02 = 47.5$ give the same answer.
\end{example}

\section{Symmetry arguments}

\begin{definition}[Symmetry argument]\label{def:iv:probability-i:symmetry}
A \emph{symmetry argument}\index{symmetry argument} computes a probability or an expectation by exhibiting a
transformation of the sample space that preserves probabilities and exchanges the outcomes of interest, so that
they must have equal probability, without computing either.
\end{definition}

\begin{example}[The first ace]\label{ex:iv:probability-i:ace}
Where, on average, is the first ace in a shuffled deck? The four aces cut the 48 other cards into five gaps, and
by symmetry each gap holds on average $48/5 = 9.6$ cards, so the first ace is at position $9.6 + 1 = 53/5 = 10.6$
on average. An exhaustive count over small decks confirms the formula $(n+1)/(k+1)$.
\end{example}

The symmetry must be real. ``Each of the six strategies is equally likely to be the best next year'' is a
symmetry only if the strategies are exchangeable; a candidate who says so should say that it is an assumption.

\section{Paradoxes and what they test}

Most probability paradoxes are questions in which the answer depends on a protocol the question does not state:
how the information was generated. The interviewer who asks the hook's question wants to hear that.

\begin{proposition}[Two dice, two protocols]\label{prop:iv:probability-i:protocols}
Two fair dice are rolled. (A) If the informant reports ``at least one six'' exactly when there is one, the
chance of two sixes given the report is $\tfrac1{11}$. (B) If the informant looks at one die chosen at random
and reports that it shows six, the chance is $\tfrac16$.
\end{proposition}

\begin{proof}
(A) Eleven of the 36 equally likely outcomes contain a six, one of them two. (B) The report depends only on the
die looked at, which is independent of the other die; the other die is a six with probability $\tfrac16$
(\cref{fig:iv:probability-i:tree}).
\end{proof}

\begin{omfigure}
\begin{tikzpicture}[font=\footnotesize,
  nd/.style={draw=omDef,rounded corners=2pt,fill=omDef!6,inner sep=3pt,align=center},
  hi/.style={draw=omThm,rounded corners=2pt,fill=omThm!8,inner sep=3pt,align=center}]
\node[nd] (a0) at (0,0) {two dice\\36 outcomes};
\node[nd] (a1) at (3.1,0.9) {at least one six\\11 outcomes};
\node[nd] (a2) at (3.1,-0.9) {no six\\25 outcomes};
\node[hi] (a3) at (6.4,0.9) {both six: 1 of 11};
\draw[-{Stealth},gray] (a0) -- (a1); \draw[-{Stealth},gray] (a0) -- (a2); \draw[-{Stealth},gray] (a1) -- (a3);
\node[font=\footnotesize\bfseries,anchor=west] at (-0.9,1.75) {(A) report ``at least one six'' when true};
\node[nd] (b0) at (0,-3.9) {two dice\\36 outcomes};
\node[nd] (b1) at (3.1,-3.0) {look at a random die\\it shows six: $\tfrac16$};
\node[nd] (b2) at (3.1,-4.8) {it shows another\\value: $\tfrac56$};
\node[hi] (b3) at (6.4,-3.0) {other die six: $\tfrac16$};
\draw[-{Stealth},gray] (b0) -- (b1); \draw[-{Stealth},gray] (b0) -- (b2); \draw[-{Stealth},gray] (b1) -- (b3);
\node[font=\footnotesize\bfseries,anchor=west] at (-0.9,-2.05) {(B) report the value of a die chosen at random};
\end{tikzpicture}

\omcaption{The same sentence, ``there is a six'', under two protocols. Under (A) the report selects the eleven
outcomes with a six; under (B) it tells you about one die only, and the other die is untouched. Exact
enumeration and a simulation in \texttt{iv\_prob1.py}.}\label{fig:iv:probability-i:tree}
\end{omfigure}

\begin{method}[Answering a paradox]\label{met:iv:probability-i:paradox}
\begin{enumerate}
\item Ask, or state, the protocol: who chose what to reveal, and how.
\item Write the joint sample space of the hidden state and the report.
\item Condition on the report, not on the sentence's literal content.
\item Give the answer for each plausible protocol if the interviewer declines to choose.
\end{enumerate}
\end{method}

The classic problems of this family, Bertrand's box (1889), the two-children problem (discussed by Gardner in
1959 and analysed as a protocol question by Bar-Hillel and Falk in 1982) and the three-door problem (Selvin's
letter of 1975), appear in the bank only as variants with their own answers.

\section{Worked answers}

\begin{example}[Which venue filled it?]\label{ex:iv:probability-i:venue}
\emph{``A router sends an order to venue A with probability 0.7, where it fills with probability 0.4, and otherwise to
venue B, where it fills with probability 0.9. The order filled. What is the chance it went to A?''} Count in a
population of 100 orders: 70 go to A and 28 fill there; 30 go to B and 27 fill there. Of the 55 fills, 28 came from A:
$28/55 \approx 0.51$. Before the fill the answer was 0.7; the fill is evidence for B, which fills more often, and the
posterior falls to about a half. \emph{Check:} if both venues filled at the same rate the fill would carry no
information and the answer would stay 0.7; with the counting table the check takes one line.
\end{example}

\begin{example}[Symmetry before summation]\label{ex:iv:probability-i:order}
\emph{``Ten days of P\&L, independent and identically distributed with a continuous distribution. What is the chance
that the best day comes after the worst day? That the best day is the last day? That the last day is the best of the
ten and the first day the worst?''} Reversing the order of the days maps ``best after worst'' to ``best before worst''
and leaves the distribution unchanged, so the chance is $1/2$. All ten days are equally likely to be the best: $1/10$.
For the last, there are $10 \times 9$ equally likely (best, worst) pairs of days: $1/90$. The interviewer is scoring
whether the candidate reaches for exchangeability before writing a sum; the answer takes fifteen seconds that way and
several minutes the other.
\end{example}

\section{Question bank}

\begin{interviewq}[$\star$ \iqroles{trader, developer} \iqfirm{any}]\label{iq:iv:probability-i:1}
In how many orders can five different trades be booked? In how many distinct orders can the letters of AABBC be
arranged?
\end{interviewq}

\begin{interviewq}[$\star$ \iqroles{trader} \iqfirm{market maker}]\label{iq:iv:probability-i:2}
What is the chance of at least one six in four rolls of a fair die? Is a bet on it at even money favourable?
\end{interviewq}

\begin{interviewq}[$\star$ \iqroles{trader} \iqfirm{market maker}]\label{iq:iv:probability-i:3}
Two cards are dealt from a full deck. What is the chance both are aces?
\end{interviewq}

\begin{interviewq}[$\star$ \iqroles{developer, researcher} \iqfirm{proprietary firm}]\label{iq:iv:probability-i:4}
In how many ways can ten identical child orders be routed to four venues? In how many if each venue must
receive at least one?
\end{interviewq}

\begin{interviewq}[$\star$ \iqroles{researcher, mle} \iqfirm{systematic fund}]\label{iq:iv:probability-i:5}
Six strategies are exchangeable. What is the chance that strategy A beats strategy B next year? That A is the
best of the six? That A beats both B and C?
\end{interviewq}

\begin{interviewq}[$\star\star$ \iqroles{developer} \iqfirm{proprietary firm}]\label{iq:iv:probability-i:6}
A system draws 1\,000 client order identifiers uniformly from a million values. What is the chance that two
coincide? About how many identifiers can it draw before the chance reaches one half?
\end{interviewq}

\begin{interviewq}[$\star\star$ \iqroles{risk, researcher} \iqfirm{bank}]\label{iq:iv:probability-i:7}
One account in 500 is fraudulent. An alert fires on 95\% of fraudulent accounts and 2\% of legitimate ones. An
alert fires on an account; what is the chance it is fraudulent?
\end{interviewq}

\begin{interviewq}[$\star\star$ \iqroles{trader, researcher} \iqfirm{market maker}]\label{iq:iv:probability-i:8}
``I rolled two dice; at least one is a six. What is the chance both are?'' Give the answer under two protocols,
and say which question you would ask the interviewer.
\end{interviewq}

\begin{interviewq}[$\star\star$ \iqroles{trader} \iqfirm{proprietary firm}]\label{iq:iv:probability-i:9}
On average, how many cards do you turn over from a shuffled deck to see the first ace? Why does the argument
not need any summation?
\end{interviewq}

\begin{interviewq}[$\star\star$ \iqroles{risk, mle} \iqfirm{bank}]\label{iq:iv:probability-i:10}
In \cref{iq:iv:probability-i:7}, a second, independent alert with the same rates also fires on the account.
What is the chance now?
\end{interviewq}

\begin{interviewq}[$\star\star\star$ \iqroles{trader, researcher} \iqfirm{market maker}]\label{iq:iv:probability-i:11}
Four envelopes contain tickets: two winning; one winning and one losing; two losing; two winning and one losing.
You choose an envelope at random and draw a ticket at random: it is winning. You draw a second ticket from the
same envelope. What is the chance it is winning too?
\end{interviewq}

\begin{interviewq}[$\star\star\star$ \iqroles{trader} \iqfirm{proprietary firm}]\label{iq:iv:probability-i:12}
Four doors, one prize. You pick a door; the host, who knows where the prize is, opens one of the other three
that is empty, choosing at random among the empty ones. You may stay, or switch to one of the two other closed
doors, chosen at random. What is your chance of winning each way?
\end{interviewq}

\begin{interviewq}[$\star\star\star$ \iqroles{developer, researcher} \iqfirm{any}]\label{iq:iv:probability-i:13}
Two orders arrive at independent uniform times in the same minute. What is the chance they arrive within ten
seconds of each other?
\end{interviewq}

\begin{interviewq}[$\star\star\star$ \iqroles{researcher, trader} \iqfirm{systematic fund}]\label{iq:iv:probability-i:14}
A trader has two positions, each profitable with probability one half and each opened on a day of the week
chosen uniformly, independently. You learn that at least one of them is profitable and was opened on a Monday.
What is the chance that both are profitable? Why is the answer not one third?
\end{interviewq}

\begin{omsources}
\item J.~Bertrand, \emph{Calcul des probabilit\'es}, Gauthier-Villars, 1889 (the box paradox).
\item M.~Gardner, ``Mathematical games'', \emph{Scientific American}, May 1959 (the two-children problem).
\item M.~Bar-Hillel and R.~Falk, ``Some teasers concerning conditional probabilities'', \emph{Cognition} 11(2),
1982, 109--122.
\item S.~Selvin, letter to the editor, \emph{The American Statistician} 29(1), 1975, 67.
\item One Quant Book~4, chapters 1 and~14 (conditional expectation; Bayesian methods).
\end{omsources}
